\section{Ascensor}
\textit{Dificultate: 1/5} \\
\textit{Propunător: Matteo-Alexandru Verzotti} \\
\textit{Developer: Matteo-Alexandru Verzotti} \\
\textit{First solve: IvanAndrei după 00:05:11}  \\

Putem separa studenții care urcă și cei care coboară în două grupe diferite. La urcare, liftul va prelua doar studenții care urcă, iar la coborâre, doar studenții care coboară.

Pentru fiecare din cele două grupe putem ține o listă \textit{urcă}, respectiv \textit{coboară}.
Fiecare student care fie urcă, fie coboară, va marca în lista corespunzătoare a sa pe pozițiile $s$ și $t$ valoarea $1$ -- un indicator că se va face o oprire la etajele respective.

În final, numărul de opriri va fi egal cu numărul de valori egale cu $1$ din cele două liste. Pentru $100$ de puncte, trebuie avut grijă la etajul $N - 1$; dacă există și urcări și coborâri la etajul respectiv, acestea trebuie numărate o singură dată.

\medskip

Complexitate finală: $\mathcal{O}(N)$.

